\documentclass[12pt,reqno]{amsart}
%%%%%%%%%%%%%%%%%%%%%%%%%%%%%%%%%%%%%%%%%%%%%%%%%%%
%								Packages									         %
%%%%%%%%%%%%%%%%%%%%%%%%%%%%%%%%%%%%%%%%%%%%%%%%%%%
\usepackage[T1]{fontenc}
\usepackage{amsmath, amssymb, amsthm, amscd, amsfonts}
\usepackage{mathtools}
\usepackage{enumerate,multicol}
\usepackage[shortlabels]{enumitem}

\usepackage[dvipsnames]{xcolor}
\definecolor{DartmouthGreen}{HTML}{00693E}
\definecolor{DartmouthDarkGreen}{HTML}{004B23}
\definecolor{DartmouthLightGreen}{HTML}{4BAF7A}

\usepackage[
  colorlinks=true,
  hyperindex,
  linkcolor=DartmouthDarkGreen,
  citecolor=DartmouthGreen,
  urlcolor=DartmouthGreen,
  pagebackref=true
]{hyperref}

\usepackage{parskip}
\usepackage[letterpaper, margin=.5in, top=1in, bottom=1.4in, headheight=42pt, headsep=14pt, footskip=50pt]{geometry}
\linespread{1.1}

\usepackage{algorithm,algpseudocode,verbatim}
\usepackage{float,caption,comment}

\usepackage{tikz,tikz-cd}
\usetikzlibrary{graphs, graphs.standard,positioning}



\usepackage{calrsfs}
\DeclareMathAlphabet{\pazocal}{OMS}{zplm}{m}{n}
\usepackage{newcent,fouriernc}

%%%%%%%%%%%%%%%%%%%%%%%%%%%%%%%%%%%%%%%%%%%%%%%%%%%
%								Copyright 								         %
%%%%%%%%%%%%%%%%%%%%%%%%%%%%%%%%%%%%%%%%%%%%%%%%%%%
\usepackage{ccicons}
\usepackage{fancyhdr}

\pagestyle{fancy}
\fancyhf{}
\fancyfoot[R]{\footnotesize \thepage}
\renewcommand{\headrulewidth}{0pt}
\renewcommand{\footrulewidth}{0pt}

\newcommand{\originalurl}{}
\newcommand{\originalurlI}{}

\fancypagestyle{firstpage}{%
  \fancyhf{}
\fancyhead[L]{\footnotesize Dartmouth College \\ Math 121: Commutative Algebra}
\fancyhead[R]{\footnotesize \href{https://www.juliettebruce.xyz/teaching/121S26}{Spring 2026}}
\fancyfoot[L]{\footnotesize \ccbyncsa\ \ \ccCopy \ Juliette Bruce 2026.\\
Licensed under \href{https://creativecommons.org/licenses/by-nc-sa/4.0/}{Creative Commons BY-NC-SA 4.0}.\\
Adapted from \ccCopy \href{https://sites.lsa.umich.edu/kesmith/}{Karen Smith} (2019); see \href{\originalurl}{original}.}
\fancyfoot[R]{\footnotesize \thepage}
\renewcommand{\headrulewidth}{0.4pt}
\renewcommand{\footrulewidth}{0.4pt}
}

\fancypagestyle{firstpage2}{%
  \fancyhf{}
\fancyhead[L]{\footnotesize Dartmouth College \\ Math 121: Commutative Algebra}
\fancyhead[R]{\footnotesize \href{https://www.juliettebruce.xyz/teaching/121S26}{Spring 2026}}
\fancyfoot[L]{\footnotesize \ccbyncsa\ \ \ccCopy \ Juliette Bruce 2026.\\
Licensed under \href{https://creativecommons.org/licenses/by-nc-sa/4.0/}{Creative Commons BY-NC-SA 4.0}.\\
Adapted from \ccCopy \href{https://sites.lsa.umich.edu/kesmith/}{Karen Smith} (2019); see \href{\originalurl}{original} and \href{\originalurlI}{original}.}
\fancyfoot[R]{\footnotesize \thepage}
\renewcommand{\headrulewidth}{0.4pt}
\renewcommand{\footrulewidth}{0.4pt}
}


\fancypagestyle{firstpageIP}{%
  \fancyhf{}
\fancyhead[L]{\footnotesize Dartmouth College \\ Math 121: Commutative Algebra}
\fancyhead[R]{\footnotesize \href{https://www.juliettebruce.xyz/teaching/121S26}{Spring 2026}}
\fancyfoot[L]{\footnotesize \ccbyncsa\ \ \ccCopy \ Juliette Bruce 2026.\\
Licensed under \href{https://creativecommons.org/licenses/by-nc-sa/4.0/}{Creative Commons BY-NC-SA 4.0}.\\
Inspired by \ccCopy\ \href{https://sites.lsa.umich.edu/kesmith/}{Karen Smith} (2019); \href{https://sites.lsa.umich.edu/kesmith/teaching/math-614-commutative-algebra/}{see}.}
\fancyfoot[R]{\footnotesize \thepage}
\renewcommand{\headrulewidth}{0.4pt}
\renewcommand{\footrulewidth}{0.4pt}
}
%%%%%%%%%%%%%%%%%%%%%%%%%%%%%%%%%%%%%%%%%%%%%%%%%%%
%								Theorems 								         %
%%%%%%%%%%%%%%%%%%%%%%%%%%%%%%%%%%%%%%%%%%%%%%%%%%%

\newtheorem{maintheorem}{Theorem}
\newtheorem{theorem}{Theorem}
\newtheorem{lemma}[theorem]{Lemma}

\newtheorem{corollary}[theorem]{Corollary}
\newtheorem{proposition}[theorem]{Proposition}
\newtheorem{conjecture}[theorem]{Conjecture}


\theoremstyle{definition}
\newtheorem{maindefinition}[maintheorem]{Definition}
\newtheorem{definition}[theorem]{Definition}
\newtheorem{setup}[theorem]{Set-Up}
\newtheorem{notation}[theorem]{Notation}
\newtheorem{convention}[theorem]{Convention}
\newtheorem{construction}[theorem]{Construction}





\setcounter{MaxMatrixCols}{18}

%%%%%%%%%%%%%%%%%%%%%%%%%%%%%%%%%%%%%%%%%%%%%%%%%%%
%								Letters									         %
%%%%%%%%%%%%%%%%%%%%%%%%%%%%%%%%%%%%%%%%%%%%%%%%%%%


\newcommand{\CC}{\mathbb{C}}
\newcommand{\FF}{\mathbb{F}}
\newcommand{\EE}{\mathbb{E}}

\newcommand{\GG}{\mathbb{G}}
\newcommand{\HH}{\mathbb{H}}
\newcommand{\II}{\mathbb{I}}
\newcommand{\KK}{\mathbb{K}}

\newcommand{\MM}{\mathbb{M}}
\newcommand{\NN}{\mathbb{N}}
\newcommand{\PP}{\mathbb{P}}
\newcommand{\QQ}{\mathbb{Q}}
\newcommand{\RR}{\mathbb{R}}
\renewcommand{\SS}{\mathbb{S}}
\newcommand{\TT}{\mathbb{T}}
\newcommand{\VV}{\mathbb{V}}
\newcommand{\WW}{\mathbb{W}}

\newcommand{\ZZ}{\mathbb{Z}}


\newcommand{\cA}{\mathcal{A}}
\newcommand{\cB}{\mathcal{B}}
\newcommand{\cC}{\mathcal{C}}
\newcommand{\cE}{\mathcal{E}}

\newcommand{\cF}{\mathcal{F}}
\newcommand{\cG}{\mathcal{G}}

\newcommand{\cH}{\mathcal{H}}
\newcommand{\cI}{\mathcal{I}}
\newcommand{\cL}{\mathcal{L}}
\newcommand{\cM}{\mathcal{M}}
\newcommand{\cO}{\mathcal{O}}
\newcommand{\cP}{\mathcal{P}}
\newcommand{\cQ}{\mathcal{Q}}
\newcommand{\cS}{\mathcal{S}}
\newcommand{\cT}{\mathcal{T}}
\newcommand{\cR}{\mathcal{R}}
\newcommand{\cU}{\mathcal{U}}
\newcommand{\cV}{\mathcal{V}}

\newcommand{\pB}{\pazocal{B}}
\newcommand{\pC}{\pazocal{C}}
\newcommand{\pO}{\pazocal{O}}
\newcommand{\pG}{\pazocal{G}}

\newcommand{\sM}{\mathsf{M}}
\newcommand{\sN}{\mathsf{N}}

\newcommand{\sQ}{\mathsf{Q}}
\newcommand{\sP}{\mathsf{P}}

\newcommand{\fm}{\mathfrak{m}}
\newcommand{\fp}{\mathfrak{p}}
\newcommand{\fq}{\mathfrak{q}}

%%%%%%%%%%%%%%%%%%%%%%%%%%%%%%%%%%%%%%%%%%%%%%%%%%%
%								Operators									         %
%%%%%%%%%%%%%%%%%%%%%%%%%%%%%%%%%%%%%%%%%%%%%%%%%%%
\DeclareMathOperator{\Id}{Id}
\DeclareMathOperator{\ev}{ev}
\DeclareMathOperator{\ord}{ord}
%
\DeclareMathOperator{\Hom}{Hom}
\DeclareMathOperator{\End}{End}
\DeclareMathOperator{\coker}{coker}
\DeclareMathOperator{\Ann}{Ann}
\DeclareMathOperator{\img}{img}
%
\DeclareMathOperator{\Frac}{Frac}
\DeclareMathOperator{\Nil}{Nil}
\DeclareMathOperator{\red}{red}
%
\DeclareMathOperator{\Spec}{Spec}
\DeclareMathOperator{\mSpec}{mSpec}
%
\DeclareMathOperator{\SL}{SL}
%
\DeclareMathOperator{\init}{in}
\DeclareMathOperator{\lex}{lex}
\DeclareMathOperator{\grlex}{grlex}
\DeclareMathOperator{\grevlex}{grevlex}
\DeclareMathOperator{\revlex}{revlex}
\DeclareMathOperator{\LT}{LT}
\DeclareMathOperator{\LM}{LM}
\DeclareMathOperator{\LC}{LC}


%%%%%%%%%%%%%%%%%%%%%%%%%%%%%%%%%%%%%%%%%%%%%%%%%%%
%								Categories									         %
%%%%%%%%%%%%%%%%%%%%%%%%%%%%%%%%%%%%%%%%%%%%%%%%%%%

\DeclareMathOperator{\comRing}{\textbf{comRing}}
\DeclareMathOperator{\Top}{\textbf{Top}}
\DeclareMathOperator{\Set}{\textbf{Set}}
\DeclareMathOperator{\alg}{\textbf{alg}}
\DeclareMathOperator{\Mod}{\textbf{Mod}}


%%%%%%%%%%%%%%%%%%%%%%%%%%%%%%%%%%%%%%%%%%%%%%%%%%%
%								MISC									         %
%%%%%%%%%%%%%%%%%%%%%%%%%%%%%%%%%%%%%%%%%%%%%%%%%%%
\makeatletter
\newcommand*{\tp}{%
  {\mathpalette\@transpose{}}%
}
\newcommand*{\@transpose}[2]{%
  % #1: math style
  % #2: unused
  \raisebox{\depth}{$\m@th#1\intercal$}%
}
\makeatother

\newcommand{\juliette}[1]{{\color{purple} \sf  Juliette: [#1]}}
%%%%%%%%%%%%%%%%%%%%%%%%%%%%%%%%%%%%%%%%%%%%%%%%%%%
%%%%%%%%%%%%%%%%%%%%%%%%%%%%%%%%%%%%%%%%%%%%%%%%%%%
%\renewcommand{\originalurl}{https://sites.lsa.umich.edu/kesmith/wp-content/uploads/sites/1309/2024/07/NullstellensatzWorksheet.pdf}
\title{Worksheet 4.1: Modules Pt. I}

%%%%%%%%%%%%%%%%%%%%%%%%%%%%%%%%%%%%%%%%%%%%%%%%%%%
%%%%%%%%%%%%%%%%%%%%%%%%%%%%%%%%%%%%%%%%%%%%%%%%%%%

%%%%%%%%%%%%%%%%%%%%%%%%%%%%%%%%%%%%%%%%%%%%%%%%%%%
%%%%%%%%%%%%%%%%%%%%%%%%%%%%%%%%%%%%%%%%%%%%%%%%%%%

\begin{document}

%%%%%%%%%%%%%%%%%%%%%%%%%%%%%%%%%%%%%%%%%%%%%%%%%%%
%%%%%%%%%%%%%%%%%%%%%%%%%%%%%%%%%%%%%%%%%%%%%%%%%%%

\maketitle
\thispagestyle{firstpageIP}

Throughout this course, ``ring'' means \textit{commutative} ring with unity. In this worksheet $\KK$ will denote a field. A \emph{$R$-module} is an abelian group $(M, +)$ together with a scalar multiplication $R\times M \to M$, $(r,m) \mapsto rm$, satisfying the following:
\begin{enumerate}[(i)]
	\item Associativity: $(rs)m = r(sm)$ for all $r,s \in R$ and $m \in M$,
	\item Distribution over $R$: $(r+s)m = rm + sm$ for all $r,s\in R$ and $m\in M$, 
	\item Distribution over $M$: $r(m+n) = rm+rn$ for all $r\in R$ and $m,n\in M$, and
	\item Identity: $1_{R}m = m$ for all $m\in M$. 
\end{enumerate}
Since $M$ is an abelian group, it has a distinguished additive identity, which we normally denote by $0$ or by $0_{M}$ if we need to distinguish it from $0_{R}$. When the base ring is clear from context, we might simply say that $M$ is a module, or we might say that $M$ is a module over $R$. As we will see, modules are the analogue of vector spaces in which the scalars come from a ring instead of a field. They include ideals, quotient rings, abelian groups, vector spaces, and many other naturally occurring algebraic objects.  

If $M$ is an $R$-module, a subset $N \subseteq M$ is an \emph{$R$-submodule} of $M$ if and only if $rm+sn \in N$ for all $r,s\in R$ and $m,n\in N$. That is, a submodule is a subgroup of $M$ such that the scalar multiplication of $R$ on $M$ restricts to a scalar multiplication on $N$. 

\hrulefill

\begin{enumerate}[leftmargin=*, series=problems]
\item \textbf{Basic Identities.} Let $R$ be a ring and $M$ an $R$-module. Prove the following identities.
	\begin{enumerate}
	\item $r0_{M} = 0_{M}$ for all $r \in R$.
	\item $0_{R}m = 0_{M}$ for all $m \in M$. 
	\item $(-r)m=r(-m)$ for all $r \in R$ and all $m\in M$. 
	\end{enumerate}
	
\item \textbf{Basic Examples of Modules.} 
	\begin{enumerate}
	\item If $\KK$ is a field, explain why a $\KK$-module is the same as a $\KK$-vector space.
	\item Explain why $R$ is itself an $R$-module in a natural way.
	\item Show that if $I \subseteq R$ is an ideal, then the multiplication on $R$ naturally makes $I$ into an $R$-module.
	\item Prove that if $I \subseteq R$ is an ideal, then $R/I$ is an $R$-module.
	\item Prove that a $\ZZ$-module is the same thing as an abelian group.
	\item If $I \subseteq R$ is an ideal and $M$ is an $R$-module, show that
\[
IM\coloneqq \left\{\sum_{i=1}^{n} r_i m_i \;\middle|\; n\geq 0,\ r_i\in I,\ m_i\in M\right\}
\]
is a submodule of $M$.
	\end{enumerate}

\item \textbf{Submodules Generated by Elements.} Let $M$ be an $R$-module. Given elements $m_{1},\ldots,m_{t} \in M$, the submodule of $M$ generated by $m_{1},\ldots,m_{t}$ is the set
\[
Rm_{1}+Rm_{2}+\cdots+Rm_{t}\coloneqq \left\{r_1m_1+\cdots+r_tm_t \;\middle|\; r_1,\ldots,r_t \in R \right\} \subseteq M. 
\]
We say that $m_{1},\ldots,m_{t}$ are \emph{generators} for $M$ if $Rm_{1}+Rm_{2}+\cdots+Rm_{t}=M$.  A submodule $N \subseteq M$ is a \emph{cyclic submodule} if $N=Rm$ for some $m\in M$. 
	\begin{enumerate}
	\item Prove that $Rm_{1}+Rm_{2}+\cdots+Rm_{t}$ is an $R$-submodule of $M$.
	\item Prove that $Rm_{1}+Rm_{2}+\cdots+Rm_{t}$ is the smallest $R$-submodule of $M$ containing $m_{1},\ldots,m_{t}$. 
	\item Prove that a ring $R$ is itself a cyclic $R$-module by finding a generator. 
	\item In the $\ZZ$-module $\ZZ/12\ZZ$, compute the submodules generated by $\overline{8}$ and $\overline{6}$.
	\item Let $R=\KK[x]$ and $M=R/\langle x^3 \rangle$. Describe the submodules generated by $\overline{1}$, $\overline{x}$, and $\overline{x}^2$. 
	\end{enumerate}
\end{enumerate}

\hrulefill

A \emph{homomorphism} of $R$-modules is a group homomorphism $\phi:M \to N$ such that $\phi(rm)=r\phi(m)$ for all $r\in R$ and $m\in M$. Two $R$-modules $M$ and $N$ are \emph{isomorphic} if and only if there exist $R$-module homomorphisms $\phi:M \to N$ and $\psi:N\to M$ such that $\phi\circ \psi = \Id_{N}$ and $\psi\circ \phi=\Id_{M}$. Equivalently, there is a bijective homomorphism of $R$-modules $\phi:M\to N$. We call such maps \emph{isomorphisms} of $R$-modules.  If $\phi:M \to N$ is a homomorphism of $R$-modules, there are two natural submodules associated to it: its \emph{kernel} and its \emph{image}:
\[
\ker(\phi) \coloneqq \left\{ m \in M \;\; \big| \;\; \phi(m)=0_{N} \right\}\subseteq M \quad \quad \text{and} \quad \quad \img(\phi) \coloneqq \left\{ \phi(m) \;\; \big| \;\; m \in M\right\} \subseteq N.
\]
One should readily check that both the kernel and image of a module homomorphism are submodules. As with linear maps, group homomorphisms, and ring homomorphisms, the kernel measures the failure of injectivity and the image measures the failure of surjectivity.

Given a submodule $N \subseteq M$, the quotient abelian group $M/N$ inherits an $R$-module structure via
\[
  r(m+N) \coloneqq rm + N.
\]
Well-definedness uses that $N$ is a submodule, not merely a subgroup. The projection $\pi \colon M \twoheadrightarrow M/N$ is a homomorphism with $\ker(\pi) = N$. The quotient $M/N$ is the module obtained by forcing $N = 0$.  Precisely: if $f \colon M \to P$ is a homomorphism with $N \subseteq \ker(f)$, there is a unique $\overline{f} \colon M/N \to P$ with $f = \overline{f} \circ \pi$:
\[
\begin{tikzcd}[row sep = 3.5em, column sep = 3.5em]
  M \arrow[r,"f"] \arrow[d,two heads,"\pi"'] & P \\
  M/N \arrow[ur,dashed,"\overline{f}"'] &
\end{tikzcd}
\]
Phrased differently, there is a natural isomorphism 
\[
\Hom_R(M/N,\, P) \cong \left\{ f \in \Hom_R(M,P) \mid N \subseteq \ker(f) \right\};
\] 
That is, homomorphisms out of $M/N$ are exactly homomorphisms out of $M$ that kill~$N$. This is usually summarized as the First Isomorphism Theorem. 
  
\begin{theorem}[First Isomorphism Theorem]
  Every homomorphism $\phi \colon M \to N$ of $R$-modules induces an isomorphism
  \[
    M/\ker(\phi) \xrightarrow{\ \sim\ } \img(\phi), \quad m + \ker(\phi) \longmapsto \phi(m).
  \]
\end{theorem}
 
In other words, quotienting by the kernel accounts for all the redundancy: every surjection is a quotient map.
  
\begin{theorem}[Third Isomorphism Theorem]
  Let $N \subseteq M$ be a submodule and $\pi \colon M \twoheadrightarrow M/N$ the projection.  There is an inclusion-preserving bijection
  \[
    \left\{\, \text{submodules of } M \text{ containing } N \,\right\}
    \;\xleftrightarrow{\;1:1\;}\;
    \left\{\, \text{submodules of } M/N \,\right\}
  \]
  given by $L \mapsto L/N = \pi(L)$ with inverse $\overline{L} \mapsto \pi^{-1}(\overline{L})$.  Moreover, for $N \subseteq L \subseteq M$:
  \[
    M/L \;\cong\; (M/N)\big/(L/N).
  \]
\end{theorem}
 
  The correspondence says that the submodule lattice of $M/N$ is the interval $[N, M]$ in the submodule lattice of~$M$.  The supplementary isomorphism $(M/N)/(L/N) \cong M/L$ is sometimes called the \emph{Third Isomorphism Theorem}: quotienting in two stages is the same as quotienting all at once.
 

\hrulefill

\begin{enumerate}[leftmargin=*, resume=problems]
\item \textbf{Basic Examples of Module Homomorphisms.} 
	\begin{enumerate}
		\item If $\KK$ is a field, explain why a homomorphism of $\KK$-modules is the same as a $\KK$-linear map. 
		\item Explain why a homomorphism of $\ZZ$-modules is the same as a homomorphism of abelian groups. 
	\end{enumerate} 

\item \textbf{Cyclic Modules \& Annihilators.} Let $M$ be an $R$-module. The \emph{annihilator} of an element $m\in M$ is $\Ann_{R}(m)\coloneqq \{r \in R\;\mid\; rm = 0 \}$. 
	\begin{enumerate}
	\item Prove that $\Ann_{R}(m)$ is an ideal of $R$.
	\item What is $\Ann_{R}(0_{M})$?
	\item Consider the multiplication by $m$ function: 
	\[
	\phi_{m}:R \to M \quad \quad \text{given by} \quad \quad r \mapsto rm.
	\]
	Prove that $\phi_{m}$ is a homomorphism of $R$-modules.
	\item Show that $\ker(\phi_m)=\Ann_{R}(m)$ and $\img(\phi_m)=Rm$. 
	\item Prove that a submodule $N \subseteq M$ is cyclic if and only if $N \cong R/I$ for some ideal $I \subseteq R$. 
	\end{enumerate}
\end{enumerate}

\hrulefill

Let $\{M_i\}_{i\in I}$ be a family of $R$-modules indexed by a set $I$. The \emph{direct product} is
\[
\prod_{i\in I}M_i=\left\{(m_i)_{i\in I}\mid m_i\in M_i\right\},
\]
with componentwise addition and scalar multiplication. The \emph{direct sum} is the submodule
\[
\bigoplus_{i\in I}M_i=\left\{(m_i)_{i\in I}\in \prod_{i\in I}M_i\mid m_i=0 \text{ for all but finitely many } i\right\}.
\]
For each $j \in I$ there is a projection map $\pi_j$ and an inclusion map $\iota_j$:
\[
\pi_j:\prod_{i\in I}M_i\to M_j,
\quad
(m_i)_{i\in I}\longmapsto m_j,
\quad \quad \text{and} \quad \quad 
\iota_j:M_j\to \bigoplus_{i\in I}M_i,
\quad
m\longmapsto (0,\ldots,0,m,0,\ldots).
\]
The product is characterized by maps to the factors, while the direct sum is characterized by maps from the summands. These are the universal properties that make products and sums so useful.

\hrulefill

\begin{enumerate}[leftmargin=*, resume=problems]
\item \textbf{Direct Sums and Products of Modules.} Let $\{M_{i}\}_{i\in I}$ be a family of $R$-modules indexed by a set $I$. 
	\begin{enumerate}
	\item Prove that both $\prod_{i\in I} M_{i}$ and $\bigoplus_{i\in I}M_{i}$ are $R$-modules. 
	\item Explain why, if $I$ is finite, $\prod_{i \in I} M_{i} = \bigoplus_{i \in I} M_i$. 
	\item Consider the $\ZZ$-modules $M_{n}=\ZZ/n\ZZ$ for $n\geq 1$. Describe explicitly an element of $\prod_{n=1}^{\infty}M_{n}$ that is not an
        element of $\bigoplus_{n=1}^{\infty}M_{n}$.
            \item Let $M_i=R$ for all $i\in \NN$. Show that the submodule of $\prod_{i\in \NN}R$ generated by the standard basis vectors $e_1,e_2,e_3,\ldots$ is $\bigoplus_{i\in \NN}R$, not the whole product, provided $R\neq 0$.
    \item For two modules $M$ and $N$, write down the inclusion maps $\iota_M,\iota_N$ and projection maps $\pi_M,\pi_N$ for $M\oplus N$. Verify the following identities
    \[
    \pi_M \circ \iota_M=\Id_M,
    \quad
    \pi_N \circ \iota_N=\Id_N,
    \quad
    \pi_M \circ \iota_N=0,
    \quad
    \pi_N \circ \iota_M=0,
\quad 
    \text{and}
    \quad 
    \iota_M\circ \pi_M+\iota_N\circ \pi_N=\Id_{M\oplus N}.
    \]

        \end{enumerate}

\item \textbf{Universal Properties of Products and Sums.} Let $\{M_{i}\}_{i\in I}$ be a family of $R$-modules indexed by $I$. 
    \begin{enumerate}
    \item Let $L$ be an $R$-module and suppose that for every $i\in I$ we are given a homomorphism $\phi_i:L\to M_i$. Prove that there exists a unique homomorphism $\phi:L\to \prod_{i\in I}M_i$ such that the following diagram commutes for all $i\in I$:
    \[
    \begin{tikzcd}[row sep = 3.5em, column sep = 4.5em]
    L \arrow[r, dashed, "\phi"] \arrow[dr, swap, "\phi_i"] & \displaystyle \prod_{i \in I} M_{i} \arrow[d, "\pi_i"] \\
    & M_{i}
	\end{tikzcd}.
	\]
    \item Let $L$ be an $R$-module and suppose that for every $i\in I$ we are given a homomorphism $\psi_i:M_i\to L$. Prove that there exists a unique homomorphism $\psi:\bigoplus_{i\in I} M_{i} \to L$ such that the following diagram commutes for all $i\in I$:
    \[
    \begin{tikzcd}[row sep = 3.5em, column sep = 4.5em]
    \displaystyle \bigoplus_{i\in I} M_{i} \arrow[r, dashed, "\psi"] & L\\
    M_{i} \arrow[u, swap, "\iota_{i}"] \arrow[ur, swap, "\psi_{i}"]&
	\end{tikzcd}.
	\]
    \item Deduce natural isomorphisms of $R$-modules
    \[
    \Hom_R\left(L, \; \prod_{i\in I}M_i\right)\cong \prod_{i\in I}\Hom_R(L,\; M_i)
    \quad \quad \text{and} \quad \quad 
    \Hom_R\left(\bigoplus_{i\in I}M_i,\; L\right)\cong \prod_{i\in I}\Hom_R(M_i,\; L).
    \]
    \item Explain why a family of maps $\psi_i:M_i\to L$ does not generally define a map $\prod_{i \in I} M_i\to L$. Use the case $M_i=L=\ZZ$ and $\psi_i=\Id_{\ZZ}$ to illustrate the problem.
    \item When $I$ is finite, explain why the two universal properties describe the same module. This is why finite direct sums are also finite direct products.
    \end{enumerate}
\end{enumerate}

\hrulefill

Over a field, every module (i.e.\ vector space) has a basis, and a linear map is determined by its values on a basis. Over a general ring, modules need not have bases at all---but those that do are of fundamental importance. The goal of this section is to set up the theory of free modules and to understand precisely how they differ from vector spaces.
  
Let $M$ be an $R$-module. A subset $B \subseteq M$ \emph{generates} (or \emph{spans}) $M$ if every element of $M$ can be written as a finite $R$-linear combination of elements of $B$, i.e.\ if
\[
M = \left\{\sum_{i=1}^n r_i m_i \;\bigg|\; n\geq 0,\; r_i\in R,\; m_i\in B\right\}.
\]
We say $B$ is \emph{linearly independent} if, whenever $m_1,\ldots,m_n\in B$ are distinct and $\sum_{i=1}^n r_i m_i = 0$, we have $r_i=0$ for all $i$. A \emph{basis} for $M$ is a subset that is both generating and linearly independent, or equivalently, a subset $\{b_s\}_{s\in S}$ such that every element of $M$ can be written \emph{uniquely} as a finite $R$-linear combination $\sum_s r_s b_s$.
 
Over a field, these notions interact exactly as in undergraduate linear algebra: every spanning set contains a basis, every linearly independent set extends to one, and all bases have the same cardinality. Over a general ring, the first two statements can fail; in particular, many modules do not admit bases at all. A module that does admit a basis is called \emph{free}. Given a set $S$, the \emph{free $R$-module on $S$} is
\[
R^{(S)}\coloneqq \bigoplus_{s\in S}R.
\]
Concretely, $R^{(S)}$ consists of all functions $S\to R$ with finite support. It has a distinguished element $e_s$ for each $s\in S$, where $e_s$ has a $1$ in the $s$-coordinate and $0$ elsewhere. Every element of $R^{(S)}$ can then be written uniquely as a finite sum
\[
\sum_{s\in S}r_se_s,
\]
so $\{e_s\mid s\in S\}$ is a basis, called the \emph{standard basis}. When $S=\{1,\ldots,n\}$ we write $R^n$ instead of $R^{(S)}$.
 
The key feature of a free module is that a homomorphism out of it is completely determined by where the basis goes. Precisely: for every function of sets $\alpha\colon S\to M$ (where $M$ is any $R$-module), there is a unique $R$-module homomorphism $\widetilde{\alpha}\colon R^{(S)}\to M$ such that $\widetilde{\alpha}(e_s)=\alpha(s)$ for all $s\in S$. In categorical language, the functor $S\mapsto R^{(S)}$ is left adjoint to the forgetful functor from $R$-modules to sets. This universal property characterizes $R^{(S)}$ up to unique isomorphism and is the reason free modules play a role analogous to that of vector spaces---they are the modules we can map out of ``for free.''

\hrulefill

\begin{enumerate}[leftmargin=*, resume=problems]
\item \textbf{Free Modules and Their Universal Property.}
    \begin{enumerate}
    \item Verify directly that every element of $R^{(S)}$ can be written uniquely as a finite $R$-linear combination of the standard basis elements $e_s$.
    \item Prove the universal property of $R^{(S)}$ stated above.
    \item Deduce that there is a natural isomorphism
    \[
    \Hom_R(R^{(S)},M)\cong \prod_{s\in S}M.
    \]
    Why is the right side a product rather than a direct sum?
    \item Show that an $R$-module $M$ is generated by $m_1,\ldots,m_n$ if and only if there is a surjective homomorphism $R^n\twoheadrightarrow M$ sending $e_i$ to $m_i$.
    \item Conclude that every finitely generated $R$-module is a quotient of a finite free module.
    \item More generally, show that every $R$-module is a quotient of a free $R$-module.
    \item For $M=R/I$, write $M$ as a quotient of a free module and identify the kernel.
    \item For $R=\KK[x]$ and $M=R/\langle x^3\rangle$, describe $M$ as a module with one generator and one relation.
    \end{enumerate}

\item \textbf{Free Modules Are Not Quite Vector Spaces.}
    \begin{enumerate}
    \item Show that $\ZZ/n\ZZ$ is not a free $\ZZ$-module for any $n>1$.
    \item More generally, if $R$ is an integral domain and $0\neq I\subsetneq R$ is an ideal, show that $R/I$ is not a free $R$-module.
    \item In the $\ZZ$-module $\ZZ$, show that $\{2,3\}$ is a generating set but not a basis.
    \item In the $\ZZ$-module $\ZZ$, show that $\{2\}$ is linearly independent but does not generate $\ZZ$.
    \item Let $R=\ZZ/6\ZZ$. Show that the element $\overline{2}\in R$ is nonzero but not part of a basis for the free $R$-module $R$.
    \item Explain why the previous examples show that bases over rings behave differently from bases over fields.
    \end{enumerate}

\item \textbf{Matrices and Maps Between Finite Free Modules.}
    \begin{enumerate}
    \item Let $A=(a_{ij})$ be an $m\times n$ matrix with entries in $R$. Define a homomorphism
    \[
    \phi_A:R^n\to R^m,
    \quad
    e_j\longmapsto \sum_{i=1}^m a_{ij}e_i.
    \]
    Prove that every homomorphism $R^n\to R^m$ arises uniquely in this way.
    \item Under this correspondence, prove that composition of homomorphisms corresponds to multiplication of matrices.
    \item When $R=\KK$ is a field, explain how this recovers the usual matrix description of linear maps between finite-dimensional vector spaces.
    \item Let $A=\begin{pmatrix}2&0\\0&3\end{pmatrix}$ and view $A$ as a homomorphism $\ZZ^2\to \ZZ^2$. Describe $\coker(A)=\ZZ^2/\img(A)$.
    \item Let $R=\KK[x]$ and consider multiplication by $x$ as a map $R^2\to R^2$. Write its matrix with respect to the standard basis, and compute its cokernel.
    \end{enumerate}
\end{enumerate}

\hrulefill

For two $R$-modules $M$ and $N$, we write $\Hom_R(M,N)$ for the set of $R$-module homomorphisms $M\to N$. Since $R$ is commutative, $\Hom_R(M,N)$ is itself an $R$-module: addition and scalar multiplication are defined pointwise by
\[
(\phi+\psi)(m)=\phi(m)+\psi(m),
\qquad
(r\phi)(m)=r\phi(m).
\]
This construction is one of the most important ways of building new modules from old ones. There is already a great deal to understand: Hom out of a free module is easy, Hom out of a quotient records annihilators, and Hom interacts cleanly with direct sums and products through their universal properties.

\hrulefill

\begin{enumerate}[leftmargin=*, resume=problems]
\item \textbf{$\Hom_R(M,N)$ as a Module.}
    \begin{enumerate}
    \item Prove that $\Hom_R(M,N)$ is an abelian group under pointwise addition. What is the zero element?
    \item Prove that the formula $(r\phi)(m)=r\phi(m)$ makes $\Hom_R(M,N)$ into an $R$-module.
    \item Show that evaluation at $1\in R$ defines an isomorphism of $R$-modules
    \[
    \Hom_R(R,M)\cong M.
    \]
    \item More generally, prove that
    \[
    \Hom_R(R/I,M)\cong \left\{m\in M\mid Im=0\right\}.
    \]
    The isomorphism should again be given by evaluation at $\overline{1}\in R/I$.
    \item Deduce that $\Hom_R(R^n,M)\cong M^n$.
    \item Prove that $\End_R(R)\cong R$ as rings, where the ring structure on $\End_R(R)$ is given by addition and composition.
    \item If $I\subseteq R$ is an ideal, identify $\Hom_R(R/I,R)$ in terms of an annihilator ideal of $R$.
    \end{enumerate}

\item \textbf{Hom Computations.}
    \begin{enumerate}
    \item Describe $\Hom_{\ZZ}(\ZZ/m\ZZ,\ZZ/n\ZZ)$ explicitly for positive integers $m,n$. Show that it is isomorphic to $\ZZ/d\ZZ$, where $d=\gcd(m,n)$.
    \item Prove that $\Hom_{\ZZ}(\QQ,\ZZ)=0$.
    \item Prove that $\Hom_{\ZZ}(\ZZ,M)\cong M$ for every abelian group $M$.
    \item Let $R=\KK[x]$ and let $M=R/\langle f\rangle$, $N=R/\langle g\rangle$. Use the preceding computation of $\Hom_R(R/I,-)$ to describe $\Hom_R(M,N)$ as a submodule of $N$.
    \item In the special case $f=x^a$ and $g=x^b$, compute $\Hom_{\KK[x]}(\KK[x]/\langle x^a\rangle,\KK[x]/\langle x^b\rangle)$ up to isomorphism.
    \end{enumerate}

\item \textbf{Functoriality of Hom.} Let $M,M',N,N'$ be $R$-modules.
    \begin{enumerate}
    \item If $\alpha:M'\to M$ is a homomorphism, prove that precomposition defines an $R$-module homomorphism
    \[
    \alpha^*:\Hom_R(M,N)\to \Hom_R(M',N),
    \quad
    \phi\longmapsto \phi\circ \alpha.
    \]
    \item If $\beta:N\to N'$ is a homomorphism, prove that postcomposition defines an $R$-module homomorphism
    \[
    \beta_*:\Hom_R(M,N)\to \Hom_R(M,N'),
    \quad
    \phi\longmapsto \beta\circ \phi.
    \]
    \item Check that $(\Id_M)^*=\Id_{\Hom_R(M,N)}$ and $(\Id_N)_*=\Id_{\Hom_R(M,N)}$.
    \item If $M''\xrightarrow{\gamma}M'\xrightarrow{\alpha}M$, prove that $(\alpha\circ\gamma)^*=\gamma^*\circ \alpha^*$.
    \item If $N\xrightarrow{\beta}N'\xrightarrow{\delta}N''$, prove that $(\delta\circ\beta)_*=\delta_*\circ \beta_*$.
    \item Let $M^*\coloneqq \Hom_R(M,R)$ be the \emph{dual module}. If $M=R^n$, describe the dual basis of $M^*$ and prove that $M^*\cong R^n$.
    \end{enumerate}
\end{enumerate}

\hrulefill

Many natural constructions in algebra are \emph{bilinear}: matrix multiplication, the dot product, and the evaluation pairing between a module and its dual. In each case, the map is additive in each variable separately, but not additive as a function of both variables together---so it is not a homomorphism from the product module. The tensor product is the device that resolves this tension. It provides a module $M\otimes_R N$ that converts bilinear maps \emph{out of} $M\times N$ into honest linear maps out of $M\otimes_R N$.
 
Let $R$ be a commutative ring with identity, and let $M$, $N$, and $P$ be $R$-modules.  An \emph{$R$-bilinear map} from $M\times N$ to $P$ is a function
\[
B\colon M\times N\to P
\]
satisfying the following three conditions for all $r\in R$, $m,m'\in M$, and $n,n'\in N$:
\begin{enumerate}[(i)]
\item $B(m+m',\,n)=B(m,n)+B(m',n)$,
\item $B(m,\,n+n')=B(m,n)+B(m,n')$,
\item $B(rm,\,n)=rB(m,n)=B(m,\,rn)$.
\end{enumerate}
In other words, $B$ is $R$-linear in each argument when the other is held fixed. A \emph{tensor product} of $M$ and $N$ is a pair $(M\otimes_R N,\;\tau)$ consisting of an $R$-module $M\otimes_R N$ and a bilinear map
\[
\tau\colon M\times N\to M\otimes_R N,
\qquad
(m,n)\longmapsto m\otimes n,
\]
that is \emph{universal} in the following sense: for every $R$-module $P$ and every $R$-bilinear map $B\colon M\times N\to P$, there exists a \emph{unique} $R$-module homomorphism $\widetilde{B}\colon M\otimes_R N\to P$ such that $B=\widetilde{B}\circ\tau$.
\[
\begin{tikzcd}
M\times N \arrow[rr,"\tau"] \arrow[dr,"B"'] & & M\otimes_R N \arrow[dl,dashed,"\widetilde{B}"] \\
& P &
\end{tikzcd}
\]
The content of this property is that $\tau$ is the ``most general'' bilinear map out of $M\times N$: every other one factors through it, and does so in exactly one way. One important consequence is that to define a linear map \emph{out of} a tensor product, it suffices to specify its values on the generators $m\otimes n$ and verify that the assignment is bilinear---the universal property then guarantees a unique, well-defined homomorphism. This strategy is used repeatedly in the exercises below.

The universal property tells us what a tensor product must \emph{do}, but not that one exists. Here is a direct construction. Begin with the free $R$-module $F$ on the set $M\times N$; its basis consists of formal symbols $[m,n]$, one for each pair $(m,n)\in M\times N$. Let $K\subseteq F$ be the submodule generated by all elements of the forms
\begin{align*}
[m+m',\,n]-[m,n]-[m',n],\\
[m,\,n+n']-[m,n]-[m,n'],\\
[rm,\,n]-r\,[m,n],\\
[m,\,rn]-r\,[m,n],
\end{align*}
for all $r\in R$, $m,m'\in M$, and $n,n'\in N$.  Set $M\otimes_R N=F/K$, and write $m\otimes n$ for the image of $[m,n]$ in the quotient. By construction, the map $(m,n)\mapsto m\otimes n$ is bilinear, and every element of $M\otimes_R N$ is a finite $R$-linear combination of such ``simple tensors.''
 
The identities that $m\otimes n$ satisfies are therefore precisely the bilinearity relations:
\begin{align*}
(m+m')\otimes n &= m\otimes n+m'\otimes n,\\
m\otimes(n+n') &= m\otimes n+m\otimes n',\\
(rm)\otimes n &= r(m\otimes n)=m\otimes(rn).
\end{align*}
 
Not every element of $M\otimes_R N$ is a simple tensor $m\otimes n$; a general element is a finite sum $\sum_i m_i\otimes n_i$. A common source of errors is to assume that every tensor can be written as a single $m\otimes n$.
 
\hrulefill

\begin{enumerate}[leftmargin=*, resume=problems]
\item \textbf{Constructing Tensor Products.}
    \begin{enumerate}
    \item Using the generators-and-relations construction above, prove the identities
    \[
    (m+m')\otimes n=m\otimes n+m'\otimes n,
    \quad
    m\otimes(n+n')=m\otimes n+m\otimes n',
    \]
    and
    \[
    (rm)\otimes n=r(m\otimes n)=m\otimes(rn).
    \]
    \item Prove that $0\otimes n=0$, $m\otimes 0=0$, and $(-m)\otimes n=-(m\otimes n)=m\otimes(-n)$.
    \item Prove that the construction described above satisfies the universal property of the tensor product.
    \item Prove that the tensor product is unique up to unique isomorphism: if $T$ and $T'$ both satisfy the universal property for $M$ and $N$, then there is a unique isomorphism $T\cong T'$ compatible with the two bilinear maps from $M\times N$.
    \item Use the universal property to construct a natural isomorphism
    \[
    M\otimes_R N\cong N\otimes_R M.
    \]
    \item If $f:M\to M'$ and $g:N\to N'$ are homomorphisms, construct a homomorphism
    \[
    f\otimes g:M\otimes_R N\to M'\otimes_R N'
    \]
    satisfying $(f\otimes g)(m\otimes n)=f(m)\otimes g(n)$.
    \end{enumerate}

\item \textbf{Basic Tensor Product Computations.}
    \begin{enumerate}
    \item Prove that the map
    \[
    R\otimes_R M\to M,
    \quad
    r\otimes m\longmapsto rm
    \]
    is an isomorphism of $R$-modules.
    \item Prove that $R^n\otimes_R M\cong M^n$.
    \item More generally, prove that
    \[
    \left(\bigoplus_{i\in I}M_i\right)\otimes_R N\cong \bigoplus_{i\in I}(M_i\otimes_R N).
    \]
    \item Let $I\subseteq R$ be an ideal. Prove that
    \[
    (R/I)\otimes_R M\cong M/IM,
    \]
    where $IM$ is the submodule defined earlier. What is the image of $\overline{r}\otimes m$ under this isomorphism?
    \item Deduce that for ideals $I,J\subseteq R$ one has $ (R/I)\otimes_R(R/J)\cong R/(I+J)$.
    \item Compute $\ZZ/m\ZZ\otimes_{\ZZ}\ZZ/n\ZZ$ for positive integers $m,n$.
    \item Compute $\QQ\otimes_{\ZZ}\ZZ/n\ZZ$.
    \item If $V$ and $W$ are finite-dimensional vector spaces over a field $\KK$, prove that
    \[
    \dim_{\KK}(V\otimes_{\KK}W)=\dim_{\KK}(V)\dim_{\KK}(W).
    \]
    \end{enumerate}

\item \textbf{Tensor Products as Linearized Bilinear Maps.}
    \begin{enumerate}
    \item Let $e_1,\ldots,e_m$ be the standard basis of $R^m$ and $f_1,\ldots,f_n$ the standard basis of $R^n$. Show that $R^m\otimes_R R^n$ is free with basis $\{e_i\otimes f_j\mid 1\leq i\leq m,\ 1\leq j\leq n\}$.
    \item Let $B:R^m\times R^n\to P$ be an $R$-bilinear map. Prove that $B$ is uniquely determined by the $mn$ elements $B(e_i,f_j)\in P$.
    \item If $M$ is generated by $m_1,\ldots,m_a$ and $N$ is generated by $n_1,\ldots,n_b$, prove that $M\otimes_R N$ is generated by the elements $m_i\otimes n_j$.
    \item Let $R=\KK[x]$. Compute $R/\langle x^a\rangle\otimes_R R/\langle x^b\rangle$ up to isomorphism.
    \item Let $I,J\subseteq R$ be ideals. The multiplication map $I\times J\to IJ$, $(a,b)\mapsto ab$, is $R$-bilinear. Use the universal property to construct a natural surjective homomorphism
    \[
    I\otimes_R J\twoheadrightarrow IJ.
    \]
    \item Show that if $I=R$, then the map in the previous part is an isomorphism. Then let $R=\KK[\varepsilon]/\langle \varepsilon^2\rangle$ and let $I=J=\langle \overline{\varepsilon}\rangle$. Compute $I\otimes_R I$ and $IJ$, and conclude that the multiplication map $I\otimes_R J\to IJ$ need not be injective.
    \end{enumerate}
\end{enumerate}

\end{document}
